\documentclass[12pt]{article}
\usepackage[T1]{fontenc}
\usepackage[utf8]{inputenc}
\usepackage{lmodern}
\usepackage{textcomp}
\usepackage{enumitem}
\usepackage{geometry}
\geometry{margin=1in}
\begin{document}
\begin{center}
Answer Key - Religious Studies - K-12 Level Assessment 19 \
\small (Answers are ordered exactly as in the question paper.)
\end{center}

\section*{Multiple Choice}
\begin{enumerate}[label=\arabic*.]
\item \textbf{Answer:} D
\\ \textit{Options:} A) Svetambara; B) Jina; C) Agama; D) Digambara
\\ \textit{Note:} The Digambara sect of Jainism believes that women are not able to achieve full renunciation and therefore permits them to lead a life of semi-renunciation, in contrast to the Svetambara sect, which allows women to attain full spiritual liberation.
\item \textbf{Answer:} C
\\ \textit{Options:} A) Prayers on behalf of people outside of the Sikh community; B) Communal confession to cleanse the community; C) Inspiration for personal meditation throughout the day; D) Doctrinal instruction for discussion in the community
\\ \textit{Note:} The purpose of hearing the Guru's Vak in the morning is to provide spiritual inspiration that guides and enhances an individual's personal meditation practice throughout the day, fostering a deeper connection to their faith and inner self.
\item \textbf{Answer:} A
\\ \textit{Options:} A) Pure Land; B) Tendai; C) Shingon; D) Soto Zen
\\ \textit{Note:} The Pure Land school of Buddhism teaches that the only hope for salvation lies in the compassionate saving grace of Amitabha Buddha (Amida), emphasizing faith and devotion as the path to rebirth in the Pure Land.
\item \textbf{Answer:} A
\\ \textit{Options:} A) Khalsa; B) Rahit; C) Panj Kakke; D) Shabad
\\ \textit{Note:} The term "Khalsa" refers to the community of initiated Sikhs who adhere to a strict code of conduct and are bound together by their commitment to Sikh principles and values, symbolizing purity and loyalty to the faith.
\item \textbf{Answer:} A
\\ \textit{Options:} A) Andal; B) Devi; C) Ganga; D) Kali
\\ \textit{Note:} Andal is recognized as the eighth-century CE female poet revered in South Indian Vaishnavism, particularly for her devotional hymns to Vishnu, making her a central figure in various temples dedicated to this deity.
\end{enumerate}

\section*{True / False}
\begin{enumerate}[label=\arabic*.]
\setcounter{enumi}{5}
\item \textbf{Answer:} True
\\ \textit{Options:} A) True; B) False
\\ \textit{Note:} In certain religious traditions, particularly in Buddhism and Hinduism, the term "vajra" is indeed used to represent both the qualities of a diamond, symbolizing strength and indestructibility, and a thunderbolt, representing spiritual power.
\item \textbf{Answer:} True
\\ \textit{Options:} A) True; B) False
\\ \textit{Note:} The socio-political alliance between the kshatriyas and shramanas was rooted in their shared opposition to Brahminic orthodoxy because both groups sought to challenge the dominance of Brahmins in religious and social hierarchies, thereby promoting alternative paths to spiritual and political authority.
\item \textbf{Answer:} True
\\ \textit{Options:} A) True; B) False
\\ \textit{Note:} Rabi'a is considered a significant figure in Sufi mysticism, known for her profound expressions of divine love in her teachings and poetry, thus affirming the truth of the statement.
\item \textbf{Answer:} True
\\ \textit{Options:} A) True; B) False
\\ \textit{Note:} In Yiddish, the term 'Bris' specifically refers to the covenant of circumcision, which is a key religious practice in Judaism symbolizing the covenant between God and the Jewish people.
\item \textbf{Answer:} False
\\ \textit{Options:} A) True; B) False
\\ \textit{Note:} The goddess known as the Holy Mother of Mount Fairy Peach is primarily associated with Chinese folk religion and is not a central figure in the religious practices of the Joseon dynasty in Korea.
\end{enumerate}

\section*{Long Form Response}
\begin{enumerate}[label=\arabic*.]
\setcounter{enumi}{10}
\item \textbf{Answer:} The communal meal in Sikhism, known as Langar, holds great significance as it embodies the principles of equality and community within the faith. Langar is served in Gurdwaras, where people from all backgrounds, regardless of caste, religion, or social status, sit together and share a meal. This practice promotes the idea that everyone is equal in the eyes of God, breaking down social barriers and fostering a sense of unity among participants. Additionally, Langar encourages selfless service and generosity, as individuals volunteer to prepare and serve food, reinforcing the community spirit central to Sikh teachings. Overall, Langar not only nourishes the body but also strengthens the bonds of community and equality among Sikhs.
\\ \textit{Note:} correct because it emphasizes how Langar serves as a practical expression of Sikhism's core values of equality and community by bringing together individuals from diverse backgrounds to share a meal, thus reinforcing the belief that all people are equal in the eyes of God.
\item \textbf{Answer:} Indonesia, as the largest Muslim country in the world, plays a vital role in shaping the cultural, social, and religious landscape of global Islam. With over 230 million Muslims, it represents a diverse interpretation of Islam that blends traditional practices with modern influences, showcasing how Islam can coexist with various cultures. This diversity contributes to a rich tapestry of Islamic thought and practice that can inspire other Muslim communities worldwide. Additionally, Indonesia's commitment to pluralism and democracy highlights the potential for peaceful coexistence among different faiths, offering a model for interfaith dialogue. Overall, Indonesia's significance lies in its ability to influence global Islamic identity and promote a more inclusive understanding of the Muslim experience.
\\ \textit{Note:} The answer correctly emphasizes Indonesia's significant role in global Islam by highlighting its large Muslim population, cultural diversity, and the unique blend of traditional and modern practices that illustrate the adaptability of Islam across different contexts.
\end{enumerate}

\vfill
\noindent\textit{End of Answer Key}
\end{document}